% ===== 7 结论（最终评委版） =====
\section{结论}
\label{sec:r-conclusion-final}

本文面向自然堆积混凝土骨料的智能筛分任务，构建了一条从原始图像到场景级级配报告的可解释处理链。方法首先利用 ImageGrains 2.0 / Cellpose-SAM 的预训练实例感知能力恢复独立颗粒，再通过成像尺度把像素几何转换为毫米级长短轴、面积和形貌特征；在此基础上，以短轴为当前筛分等效粒径基线，并通过 $w=d^3$ 将颗粒数量统计转换为质量相关代理统计，最终输出筛级比例以及 $D_{10}/D_{50}/D_{90}$。同一实例几何同时服务于二维投影形貌分类和异常候选筛查，使所有场景级结果能够追溯到具体颗粒与运行参数。

现有三张真实自然堆积样本共检测 3858 个实例，已经形成从实例分割、毫米几何到级配、形貌和异常报告的完整闭环。三张场景在当前统一参数下表现出明显不同的预测粒径结构；更重要的是，同一实例集合按数量等权统计时 $D_{50}$ 为 5.2--7.6~mm，而采用 $d^3$ 质量代理后为 14.1--27.5~mm，三个场景均出现显著向粗端移动。这一结果直接说明：自然堆积视觉中“检测了多少颗粒”不能等价为机械筛分意义下“各粒级占多少质量”，因此质量统计口径必须在方法中被显式处理。仓库中的 27 项竞赛层合成测试进一步验证了筛分映射、加权分位数、形貌、异常与报告逻辑的一致性，历史算法阶段约 27~s/张，也说明计算开销能够嵌入 30~min 的现场操作流程。

本文同时明确限制结论范围。由于当前证据集中没有同批机械筛分逐筛称量、人工实例 mask 或三维形貌/材料标签，本文不报告标准筛分绝对准确率、实例 P/R/F1 或泥团语义准确率；当前级配统一解释为基于二维几何和 $d^3$ 权重的视觉质量代理。这一边界并不削弱系统作为竞赛工程原型的完整性，而是区分了“已经由现有数据证明的功能与行为”以及“需要外部物理真值才能证明的精度”。总体上，本文形成的核心方案是\textbf{预训练实例感知 + 可追溯物理尺度 + 显式筛分语义映射 + 质量相关统计 + 分层现场质检}，其优势在于现场训练依赖低、计算链透明、输出结构完整，并为后续进一步提高物理一致性保留了清晰接口。